\subsection{一族集合的并集与交集的定义}

设 X 是一个族（§3，第4号），I 是其指标集。为了便于对下文进行直观理解，我们将 X 称为一个集合族。若 (X, I, G) 是某个集合 E 的子集族（即其目标 G 满足关系 Y ∈ G 蕴含 Y ⊂ E），则我们将该族记为 (X\_t)\_\{t ∈ I\}(X\_t ∈ G) 或简记为 (X\_t)\_\{t ∈ I\} (§3, 第6节)；为简便起见，我们将任何以 I 为指标集的集合族记为 (X\_t)\_\{t ∈ I\}.

由于关系 \((\forall x)((\iota \in I\) 并且 \(x\in X_{\iota})\Longrightarrow (x\in X_{\iota}))\) 是真的，根据S5（第一章，§4，第2节）可知，该关系

\[
(\forall \iota) (\exists Z) (\forall x) ((\iota \in I \text { 且 } x \in X _ {\iota}) \Longrightarrow (x \in Z))
\]

是真的。根据方案S8（§1，第6节），关系 \((\exists t)(t \in I\) 并且 \(x \in X_t)\) 关于 \(x\) 是集合化的.

定义1. 设 \((\mathbf{X}_i)_{i\in I}\) 设为一个集合族（相应地，为集合E的子集族）。集合 \(\mathcal{E}_x((\exists i)(i\in I\) 并且 \(x\in \mathbf{X}_i))\) ，也就是说，所有 \(x\) 属于该族中至少一个集合的 \((\mathbf{X}_{\iota})_{\iota \in \mathbf{I}}\) ，称为该族的并集，记作 \(\bigcup \mathbf{X}_{\iota}(*)\) .

如果 \((\mathbf{X}_t)_{i\in I}\) 是集合的子集族 \(E\) ，则其并集是……的子集 \(E\) ; 注意它并不依赖于 \(E\) ，也不依赖于目标 \(\mathfrak{G}\) 该映射的 \(t\to X_t\) .

显然，如果 \(\mathbf{I} = \emptyset\) ，我们有 \(\bigcup_{i \in \mathbf{I}} \mathbf{X}_i = \emptyset\) 的集合，因为关系 \((\exists i)(i \in \mathbf{I}\) 并且 \(x \in \mathbf{X}_i)\) 此时为假。

现在假设 \(I \neq \varnothing\) 。如果 \(\alpha\) 是 I 的一个元素，关系

\[
(\forall \iota) ((\iota \in I) \Longrightarrow (x \in X _ {\iota}))
\]

蕴含 \(x \in \mathbf{X}_{\alpha}\) 。因此，根据 C52（§1，第 6 节），此关系关于 \(x\) 是集合化的。

定义 2. 设 \((\mathbf{X}_{\iota})_{\iota \in \mathbf{I}}\) 是一个集合族，其指标集 I 非空。集合 \(\mathcal{E}_x((\forall \iota)((\iota \in \mathbf{I}) \Longrightarrow (x \in \mathbf{X}_{\iota}))\) ，也就是说，所有 \(x\) 它属于该族中的每一个集合 \((\mathbf{X}_{\iota})_{\iota \in \mathbf{I}}\) ，称为该族的交集，记作 \(\bigcap_{\iota} \mathbf{X}_{\iota}\) .

如果 \(I = \varnothing\) ，关系 \((\forall t)((t \in I) \Longrightarrow (x \in X_t))\) 并未在 \(x\) ；因为这是一个真实的关系，且不存在这样的集合。 \(Y\) 使得 \(x \in Y\) 是一个真关系，因为 \(Y\) 则将是所有对象的集合（参见§1，第7号，注记）。

如果 \((\mathbf{X}_t)_{t \in \mathbf{I}}\) 是集合的子集族 \(E\) 并且如果 \(I \neq \emptyset\) ，关系“ \(x \in E\) 并且 \((\forall t)((t \in I) \Longrightarrow (x \in X_t))\) 等价于

\[
(\forall \iota) ((\iota \in I) \Longrightarrow (x \in X _ {\iota}));
\]

因此它在 \(x\) 中是集体化的，并且所有满足此关系的 \(x\) 的集合等于 \(\bigcap_{i \in I} X_i\) 。如果 \(I = \emptyset\) ，关系“ \(x \in E\) 并且 \((\forall i)((i \in I) \Longrightarrow (x \in X_i))\) ”等价于 \(x \in E\) ；因此它在 \(x\) 中是集体化的，并且所有满足此关系的 \(x\) 满足此关系的 \(E\) 因此，我们可以给出如下定义：

定义3。设 \((\mathbf{X}_{\iota})_{\iota \in \mathbf{I}}\) 设E是一个集合，E的子集族。集合

\[
\mathcal {E} _ {x} (x \in \mathrm{E} \text {   且   } (\forall \iota) ((\iota \in \mathrm{I}) \Longrightarrow (x \in \mathrm{X} _ {\iota}))),
\]

换言之，所有 \(x\) （属于 \(\mathbf{E}\) 且属于每个集合 \(\mathbf{X}_t\) ）所成的集合，称为该族的交集，记作 \(\bigcap_{i=1}^{\infty} \mathbf{X}_t\) .

因此，对于 \((\mathbf{X}_{i})_{i\in\theta}\) 的一个子集族，我们有

\[
\bigcap_ {i \in \theta} X _ {i} = E.
\]

但对于 \((\mathbf{X}_t)_{i\in I}\) 的一个子集族，若其指标集非空，则交集 \(\bigcap_{i\in I}\mathbf{X}_t\) 既不依赖于E，也不依赖于 \(\iota \to \mathbf{X}_t\) ；这证明了在定义2和定义3中使用相同记号的合理性。

命题1. 设 \((\mathbf{X}_{\mathfrak{t}})_{\mathfrak{t}\in \mathbf{I}}\) 是一族集合，并设 \(f\) 是从集合K到I的一个映射。则

\[
\bigcup_ {x \in \mathbf {K}} X _ {f (x)} = \bigcup_ {i \in I} X _ {i},
\]

并且，如果 \(\mathbf{I} \neq \emptyset\) ,

\[
\bigcap_ {x \in K} X _ {f (x)} = \bigcap_ {i \in I} X _ {i}.
\]

让 \(x\) 为 \(\bigcap_{i \in I} X_i\) . 存在一个指标 \(i \in I\) 使得 \(x \in X_i\) 。由于 \(f\langle K \rangle = I\) ，存在一个指标 \(x \in K\) 使得 \(i = f(x)\) ，由此 \(x \in X_{f(x)}\) ，从而

\[
x \in \bigcup_ {x \in \mathbf {K}} X _ {f (x)}.
\]

反之，如果 \(x \in \bigcup_{x \in \mathbb{K}} X_{f(x)}\) , 存在 \(x \in \mathbb{K}\) 使得 \(x \in X_{f(x)}\) ，因此，由于 \(f(x) \in I\) , \(x \in \bigcup_{i' \in I} X_i\) 。因此

\[
\bigcup_ {x \in \mathbb {K}} X _ {f (x)} = \bigcup_ {i \in I} X _ {i}.
\]

现在假设 \(\mathbf{I} \neq \emptyset\) ，并且让 \(x\) 为 \(\bigcap_{i \in \mathbf{I}} \mathbf{X}_i\) 。对于每个元素 \(x\) 的 \(\mathbf{K}\) 我们拥有 \(f(x) \in \mathbf{I}\) ，因此 \(x \in \mathbf{X}_{f(x)}\) ，从而

\[
x \in \bigcap_ {x \in \mathbb {K}} X _ {f (x)}.
\]

反之，设 \(x\) 为 \(\bigcap_{x \in \mathbf{K}} X_{f(x)}\) 。如果 \(\iota\) 是……的任意元素 \(I\) ，存在一个元素 \(x\) 的 \(K\) 使得 \(\iota = f(x)\) ，由此 \(x \in X_t\) 并且

因此 \(x \in \bigcap_{i \in I} X_i\) 。因此

\[
\bigcap_ {x \in \mathbf {K}} X _ {f (x)} = \bigcap_ {i \in I} X _ {i}.
\]

对于给定集合的子集族，命题1的第二部分显然在去掉限制后仍然成立。 \(I \neq \emptyset\) .

推论。设 \((\mathbf{X}_t)_{i\in \mathbf{I}}\) 设为一个集合族，使得 \(\mathbf{X}_t = \mathbf{X}_x\) 对于每一对索引 \((\mathfrak{t},\mathfrak{x})\) 然后对每一个 \(\alpha \in \mathbf{I}\) 我们拥有

\[
\bigcup_ {i \in I} X _ {i} = X _ {\alpha}, \quad \text{且} (\text{若 } I \neq \varnothing) \quad \bigcap_ {i \in I} X _ {i} = X _ {\alpha}.
\]

将命题1应用于常值映射 \(\iota\to\alpha\) 的 I 到 \(\{\alpha\}\) .

定义4. 设 \(\mathfrak{F}\) 是一组集合的集合，令 \(\Phi\) 是由恒等映射定义的集合族 \(\mathfrak{F}\) 这些集合的并集 \(\Phi\) 并且（如果 \(\mathfrak{F}\) （非空时）这些集合的交集 \(\Phi\) 分别称为集合族的并集和交集 \(\mathfrak{F}\) ，并记作 \(\bigcup_{X\sim X^{\prime}}X\) 并且 \(\bigcap_{X\sim X^{\prime}}X\) .

由命题1立即得出，如果 \((\mathbf{X}_i)_{i\in I}\) 是一个集合族，那么并集以及（若 \(I\neq \emptyset\) ) 该族的交集分别等于此族元素集合的并集与交集。

